% titel: Eiskunstlaeufer: Drehimpuls und Rotationsenergie
% skript: Kap. Drehimpuls, Rotationsenergie

Beim Einziehen der Arme während einer Pirouette bleibt das System reibungsfrei:
der Drehimpuls um die vertikale Drehachse ist erhalten, obwohl das
Trägheitsmoment $I$ kleiner wird. Daraus ergeben sich die Änderungen von
Winkelgeschwindigkeit und Rotationsenergie.

\begin{werkzeugkasten}[Werkzeuge]
\textbf{Drehimpuls} einer rotierenden Anordnung um die Achse:
\begin{equation}
  L = I\,\omega\;.
\end{equation}

\textbf{Drehimpulserhaltung} (keine äusseren Drehmomente um die Achse):
\begin{equation}
  I_1\,\omega_1 = I_2\,\omega_2 = \mathrm{konst}.
\end{equation}

\textbf{Rotationsenergie}:
\begin{equation}
  E_\mathrm{rot} = \frac12\,I\,\omega^2
  = \frac{L^2}{2I}.
\end{equation}

\textbf{Hinweis:} Das Anziehen der Arme leistet Arbeit (die Armmuskeln), die
Energie zuführen kann; die Erhaltung des \emph{Drehimpulses} gilt unabhängig
davon.
\end{werkzeugkasten}

\begin{werkzeugkasten}[Vorgehen in allgemeiner Form]
Wird $I$ halbiert, so verdoppelt sich bei konstantem $L$ die
Winkelgeschwindigkeit. Für die Rotationsenergie die Formel $E_\mathrm{rot}
=\frac12 I\omega^2$ mit dem neuen Paar $(I_2,\omega_2)$ auswerten und den
Faktor bestimmen. Die Energiezunahme über die von den Armen geleistete Arbeit
erklären.
\end{werkzeugkasten}

\begin{achtung}[Typische Fallen]
\begin{itemize}
  \item Nicht $L$, sondern $I\omega$ ist konstant; die Energie ist es
       nicht.
  \item Bei $E_\mathrm{rot}=\frac12 I\omega^2$ müssen \emph{both} $I$ und
       $\omega$ geändert werden.
\end{itemize}
\end{achtung}
